\section{Special Techniques}
%==================================================

This section collects optional techniques, advanced shortcuts, and exam-safety
workflows.  Core formulas are kept in the Formula \& Key Results Reference;
this section explains when and how to use nonstandard or easily misapplied
methods.

\subsection{Limit and Algebraic Simplification Techniques}

\subsubsection{Dominant-term comparison}
\begin{extension}
Dominant-term comparison is a fast way to simplify limits of rational,
radical, exponential, or logarithmic expressions before applying a formal test.
The idea is to divide by the largest relevant scale.
\end{extension}

\begin{keyformula}[title={Dominant-Term Templates}]
\[
\frac{a_mx^m+	ext{lower powers}}{b_nx^n+	ext{lower powers}}
\sim \frac{a_m}{b_n}x^{m-n}\\quad(x\to\infty),
\]
\[
\sqrt{x^2+ax+b}=x\sqrt{1+\frac{a}{x}+\frac{b}{x^2}}
\sim x+\frac a2\quad(x\to\infty).
\]
\end{keyformula}

\begin{examtip}
At infinity, compare by growth class first:
\[
\log x \ll x^p \ll a^x \ll x! \ll x^x
\quad(p>0,\ a>1).
\]
This usually identifies whether a ratio limit is \(0\), finite nonzero, or
infinite before detailed algebra.
\end{examtip}

\begin{commonmistake}
Dominant-term comparison gives the limit only after the expression has been
put on a common scale.  Do not discard terms inside a difference such as
\(\sqrt{x^2+x}-x\); cancellation may make the next-order term decisive.
\end{commonmistake}

\begin{examplebox}[title={Radical Cancellation}]
\[
\lim_{x\to\infty}\left(\sqrt{x^2+x}-x\right)
=\lim_{x\to\infty}\frac{x}{\sqrt{x^2+x}+x}
=\frac12.
\]
\end{examplebox}

\newpage
\subsubsection{Rationalisation patterns}
\begin{extension}
Rationalisation is useful whenever a limit has a radical difference or a
hidden removable singularity.  It converts cancellation into a product or
quotient that exposes the small factor.
\end{extension}

\begin{keyformula}[title={Useful Rationalisations}]
\[
\sqrt{A}-\sqrt{B}=\frac{A-B}{\sqrt A+\sqrt B},
\qquad
\sqrt[3]{A}-\sqrt[3]{B}=\frac{A-B}{A^{2/3}+A^{1/3}B^{1/3}+B^{2/3}}.
\]
\end{keyformula}

\begin{examtip}
For limits of the form \(0/0\) with radicals, rationalise before using
L'H\^opital's rule.  The algebra is often shorter and less error-prone.
\end{examtip}

\begin{commonmistake}
When rationalising a cube-root expression, the denominator is the sum
\(A^{2/3}+A^{1/3}B^{1/3}+B^{2/3}\), not just a conjugate with two terms.
\end{commonmistake}

\newpage
\subsubsection{Logarithmic limit transformation}
\begin{extension}
For powers with a variable base and variable exponent, take logarithms first.
This turns a hard expression into a product limit.
\end{extension}

\begin{keyformula}[title={Power-Limit Log Transformation}]
If \(y(x)>0\) and \(y(x)=u(x)^{v(x)}\), then
\[
\ln y=v(x)\ln u(x),
\qquad
\lim y=\exp\!\left(\lim v(x)\ln u(x)\right)
\]
whenever the logarithmic limit exists.
\end{keyformula}

\begin{examtip}
Use this for indeterminate forms \(1^\infty\), \(0^0\), and \(\infty^0\).
Convert to an exponential only at the end.
\end{examtip}

\begin{examplebox}[title={Standard Exponential Limit}]
\[
y=\left(1+\frac{a}{x}\right)^x,
\qquad
\ln y=x\ln\left(1+\frac{a}{x}\right)\to a,
\qquad
 y\to e^a.
\]
\end{examplebox}

\newpage
\subsubsection{Squeeze theorem templates}
\begin{extension}
The squeeze theorem is especially useful when a bounded oscillatory factor is
multiplied by a term that tends to zero.
\end{extension}

\begin{keyformula}[title={Common Squeeze Forms}]
\[
|\sin u|\le 1,
\qquad
|\cos u|\le 1,
\qquad
0\le 1-\cos u\le \frac{u^2}{2},
\qquad
|\sin u|\le |u|.
\]
\end{keyformula}

\begin{examtip}
If the difficult part is bounded and the simple factor tends to zero, squeeze
rather than differentiate or expand.
\end{examtip}

\begin{examplebox}[title={Oscillatory Squeeze}]
\[
\left|x^2\sin\frac1x\right|\le x^2\to0
\quad\Rightarrow\quad
\lim_{x\to0}x^2\sin\frac1x=0.
\]
\end{examplebox}

\newpage
\subsection{Differentiation Techniques and Derivative-Based Workflows}

\subsubsection{Logarithmic differentiation}
\begin{extension}
Logarithmic differentiation is an efficient method for products, quotients,
powers with variable exponents, and functions of the form \(y=f(x)^{g(x)}\).
\end{extension}

\begin{keyformula}[title={Logarithmic Differentiation Template}]
If \(y>0\), then
\[
\frac{y'}{y}=\frac{d}{dx}\ln y.
\]
For \(y=f(x)^{g(x)}\),
\[
\ln y=g(x)\ln f(x),
\qquad
 y'=f(x)^{g(x)}\left(g'(x)\ln f(x)+g(x)\frac{f'(x)}{f(x)}\right).
\]
\end{keyformula}

\begin{commonmistake}
After differentiating \(\ln y\), multiply by \(y\).  Stopping at \(y'/y\) is
not the derivative unless the question explicitly asks for logarithmic derivative.
\end{commonmistake}

\begin{examplebox}[title={Variable Power}]
For \(y=x^x\),
\[
\ln y=x\ln x,
\qquad
\frac{y'}{y}=\ln x+1,
\qquad
 y'=x^x(\ln x+1).
\]
\end{examplebox}

\newpage
\subsubsection{Implicit differentiation workflow}
\begin{extension}
Implicit differentiation is useful when solving explicitly for \(y\) is
unpleasant or impossible.  Treat \(y\) as a function of \(x\) throughout.
\end{extension}

\begin{examtip}
Differentiate every term with respect to \(x\).  Whenever differentiating an
expression involving \(y\), multiply by \(y'\) through the chain rule.  Then
collect all \(y'\)-terms and solve.
\end{examtip}

\begin{examplebox}[title={Implicit Derivative}]
If \(x^2+xy+y^2=3\), then
\[
2x+(xy)' +2yy'=0
\quad\Rightarrow\quad
2x+y+xy'+2yy'=0,
\]
so
\[
y'=-\frac{2x+y}{x+2y}.
\]
\end{examplebox}

\begin{commonmistake}
The derivative of \(xy\) is \(x y'+y\), not just \(xy'\).  Product rule and
chain rule both matter.
\end{commonmistake}

\newpage
\subsubsection{Inverse-function derivative shortcut}
\begin{extension}
When \(f\) is one-to-one and differentiable with \(f'(a)\ne0\), the derivative
of the inverse can be computed without explicitly finding the inverse formula.
\end{extension}

\begin{keyformula}[title={Inverse Derivative Formula}]
If \(f(a)=b\), then
\[
(f^{-1})'(b)=\frac{1}{f'(a)}.
\]
Equivalently,
\[
(f^{-1})'(x)=\frac{1}{f'(f^{-1}(x))}.
\]
\end{keyformula}

\begin{examtip}
For a numerical inverse-derivative question, first solve \(f(a)=b\), then
compute \(1/f'(a)\).  Do not try to derive a full inverse unless required.
\end{examtip}

\newpage
\subsubsection{Tangent and normal line workflow}
\begin{extension}
Tangent and normal line questions are mostly bookkeeping.  The decisive step is
finding the slope at the correct point.
\end{extension}

\begin{keyformula}[title={Tangent and Normal Lines}]
At \((x_0,y_0)\),
\[
y-y_0=f'(x_0)(x-x_0)
\]
for the tangent line.  If \(f'(x_0)\ne0\), the normal slope is
\[
m_{\rm normal}=-\frac{1}{f'(x_0)}.
\]
\end{keyformula}

\begin{commonmistake}
For implicit curves, compute \(y'\) first and only then substitute the point.
Substituting too early can destroy terms needed for differentiation.
\end{commonmistake}

\newpage
\subsection{Applications of Differentiation}

\subsubsection{Optimisation endpoint checklist}
\begin{extension}
Closed-interval optimisation requires both interior critical points and
endpoints.  Many exam errors come from solving \(f'(x)=0\) correctly but
forgetting the boundary.
\end{extension}

\begin{examtip}
For optimisation on \([a,b]\):
\begin{enumerate}[leftmargin=*]
\item Build the objective as a single-variable function.
\item Determine the allowed domain from constraints.
\item Solve \(f'(x)=0\) and include points where \(f'\) does not exist.
\item Evaluate \(f\) at all critical points and endpoints.
\item Choose the maximum/minimum requested by the question.
\end{enumerate}
\end{examtip}

\begin{commonmistake}
The second derivative test classifies local behaviour.  It does not replace the
closed-interval endpoint comparison for global maxima or minima.
\end{commonmistake}

\newpage
\subsubsection{Hidden constraint extraction}
\begin{extension}
In optimisation and related-rates questions, the hard part is often not
calculus but extracting the constraint from geometry, units, or wording.
\end{extension}

\begin{examtip}
Name the variables, draw the diagram if possible, write all constraints before
substitution, and only then eliminate variables.  Keep the physical domain, such
as \(x>0\), \(0<\theta<\pi/2\), or \(V\ge0\), visible.
\end{examtip}

\begin{examplebox}[title={Box With Fixed Volume}]
For an open-top box with square base side \(x\), height \(h\), and fixed volume
\(V\), the constraint is
\[
x^2h=V\quad\Rightarrow\quad h=\frac{V}{x^2}.
\]
The surface area is
\[
S=x^2+4xh=x^2+\frac{4V}{x},
\]
so the optimisation is now one-variable with domain \(x>0\).
\end{examplebox}

\newpage
\subsubsection{Related-rates workflow}
\begin{extension}
Related-rates problems are implicit differentiation problems with time as the
independent variable.
\end{extension}

\begin{examtip}
Use the order: draw, define variables, write the geometric equation, differentiate
with respect to \(t\), then substitute numerical values at the instant of interest.
\end{examtip}

\begin{commonmistake}
Do not substitute time-specific numerical values before differentiating, unless
they are constants for all time.  For example, a ladder length is constant, but
its height and horizontal distance are not.
\end{commonmistake}

\begin{examplebox}[title={Ladder Equation}]
If a ladder of length \(L\) has horizontal distance \(x(t)\) and height \(y(t)\),
then
\[
x^2+y^2=L^2.
\]
Differentiating gives
\[
2x\frac{dx}{dt}+2y\frac{dy}{dt}=0,
\qquad
\frac{dy}{dt}=-\frac{x}{y}\frac{dx}{dt}.
\]
\end{examplebox}

\newpage
\subsubsection{Derivative sign-chart workflow}
\begin{extension}
A sign chart converts derivative information into monotonicity, extrema,
concavity, and inflection behaviour without relying on a graphing calculator.
\end{extension}

\begin{examtip}
For increasing/decreasing, test the sign of \(f'\) on intervals separated by
critical points.  For concavity, test the sign of \(f''\) on intervals separated
by candidates where \(f''=0\) or does not exist.
\end{examtip}

\begin{commonmistake}
A point where \(f''=0\) is only an inflection candidate.  Concavity must actually
change sign across the point.
\end{commonmistake}

\newpage
\subsection{Integration Setup and Standard Technique Choice}

\subsubsection{Endpoint-aware integral setup}
\begin{extension}
Before computing an integral, decide what the integral represents.  This is a
method discipline rather than a new formula, and it prevents most application
errors.
\end{extension}

\begin{examtip}
Use this setup order:
\begin{enumerate}[leftmargin=*]
\item Choose the variable of integration.
\item Mark endpoints and split at intersections, sign changes, or singularities.
\item Write the slice: top minus bottom, right minus left, shell radius times
height, cross-sectional area, force times distance, or pressure times width.
\item Simplify and integrate only after the slice is correct.
\end{enumerate}
\end{examtip}

\begin{commonmistake}
For geometric area, use nonnegative height or width.  A signed integral may
cancel positive and negative regions.
\end{commonmistake}

\newpage
\subsubsection{Substitution recognition patterns}
\begin{extension}
Substitution is best viewed as reversing the chain rule.  The key recognition
cue is an inner expression whose derivative is also present up to a constant.
\end{extension}

\begin{examtip}
Common substitution cues:
\begin{itemize}[leftmargin=*]
\item polynomial inside radical or power: try that inside expression;
\item denominator derivative appears in numerator: try denominator;
\item exponent or trig argument is nontrivial: try the argument;
\item symmetric bounds and even/odd functions: consider symmetry before substitution.
\end{itemize}
\end{examtip}

\begin{examplebox}[title={Inside-Derivative Match}]
\[
\int \frac{2x}{1+x^2}\,dx,
\qquad u=1+x^2,
\qquad du=2x\,dx,
\]
so
\[
\int \frac{2x}{1+x^2}\,dx=\int \frac{du}{u}=\ln|1+x^2|+C.
\]
\end{examplebox}

\newpage
\subsubsection{Integration by parts choice and tabular form}
\begin{extension}
Integration by parts is useful when one factor becomes simpler after repeated
differentiation.  Tabular integration by parts is a compact version for
polynomial times exponential, sine, or cosine.
\end{extension}

\begin{examtip}
Try choosing \(u\) in the order: logarithmic/inverse-trig, algebraic,
trigonometric/exponential.  If repeated differentiation of \(u\) reaches zero,
tabular integration is often efficient.
\end{examtip}

\begin{examplebox}[title={Tabular Integration by Parts}]
For \(\int x^2e^x\,dx\), differentiate \(x^2\) and repeatedly integrate \(e^x\):
\[
\begin{array}{c|c}
+\ x^2 & e^x\\
-\ 2x & e^x\\
+\ 2 & e^x\\
-\ 0 & e^x
\end{array}
\]
Thus
\[
\int x^2e^x\,dx=x^2e^x-2xe^x+2e^x+C.
\]
\end{examplebox}

\begin{commonmistake}
Tabular signs alternate.  For definite integrals, do not forget to evaluate the
resulting boundary term at both endpoints.
\end{commonmistake}

\newpage
\subsubsection{Splitting absolute-value integrals}
\begin{extension}
Absolute-value integrals are piecewise integrals.  The correct split points are
where the inside expression changes sign.
\end{extension}

\begin{keyformula}[title={Absolute-Value Split}]
If \(g(x)=0\) at \(c\in(a,b)\) and \(g\) changes sign there, then
\[
\int_a^b |g(x)|\,dx
=\int_a^c \pm g(x)\,dx+\int_c^b \pm g(x)\,dx,
\]
where the sign on each interval is chosen to make the integrand nonnegative.
\end{keyformula}

\begin{examplebox}[title={Simple Split}]
\[
\int_0^2 |x-1|\,dx
=\int_0^1(1-x)\,dx+\int_1^2(x-1)\,dx=1.
\]
\end{examplebox}

\newpage
\subsubsection{Trigonometric identity selection}
\begin{extension}
For trig integrals, the right identity depends on the parity of powers.  Do not
expand everything automatically.
\end{extension}

\begin{examtip}
For \(\sin^m x\cos^n x\), save one factor if either exponent is odd.  Use
half-angle identities mainly when both exponents are even.  For \(\tan^m x\sec^n x\),
save \(\sec^2x\) if \(n\) is even and save \(\sec x\tan x\) if \(m\) is odd.
\end{examtip}

\begin{commonmistake}
Using half-angle identities too early can turn a one-line substitution into a
long algebraic expansion.
\end{commonmistake}

\newpage
\subsection{Improper Integrals and Convergence Shortcuts}

\subsubsection{Comparison by local model}
\begin{extension}
For improper integrals, convergence is determined locally near the problematic
point: infinity, zero, a finite asymptote, or an interior singularity.
\end{extension}

\begin{examtip}
Near the problematic point, replace the integrand by its leading model.  Then
compare with a \(p\)-integral:
\[
f(x)\sim Cx^{-p}\quad(x\to\infty),
\qquad
f(x)\sim C(x-a)^{-p}\quad(x\to a^+).
\]
\end{examtip}

\begin{commonmistake}
The threshold is different at infinity and near zero:
\[
\int_1^\infty x^{-p}\,dx\text{ converges iff }p>1,
\qquad
\int_0^1x^{-p}\,dx\text{ converges iff }p<1.
\]
\end{commonmistake}

\newpage
\subsubsection{Tail estimates for improper integrals}
\begin{extension}
When an improper integral is approximated by cutting off the tail, a tail bound
is needed to justify the accuracy.
\end{extension}

\begin{examtip}
If \(0\le f(x)\le g(x)\) for \(x\ge A\) and \(\int_A^\infty g\) is easy, then
\[
0\le \int_N^\infty f(x)\,dx\le \int_N^\infty g(x)\,dx.
\]
Choose \(N\) so that the upper bound is below the required tolerance.
\end{examtip}

\begin{examplebox}[title={Power Tail}]
If \(0\le f(x)\le x^{-3}\) for \(x\ge N\), then
\[
0\le\int_N^\infty f(x)\,dx\le\int_N^\infty x^{-3}\,dx=\frac{1}{2N^2}.
\]
\end{examplebox}

\newpage
\subsection{Symmetry and Transformation Tricks for Definite Integrals}

\subsubsection{Finite-interval reflection and averaging}
\begin{extension}
Reflection pairs each point \(x\) with its mirror image \(a+b-x\) on \([a,b]\).
This is useful when direct antiderivatives are awkward but the reflected
integrand complements the original one.
\end{extension}

\begin{keyformula}[title={Reflection and Averaging}]
\[
\int_a^b f(x)\,dx=\int_a^b f(a+b-x)\,dx,
\qquad
\int_a^b f(x)\,dx=\frac12\int_a^b\bigl(f(x)+f(a+b-x)\bigr)\,dx.
\]
\end{keyformula}

\begin{examplebox}[title={Complementary Integrand}]
Let
\[
I=\int_0^1\frac{x^2}{x^2+(1-x)^2}\,dx.
\]
Reflecting \(x\mapsto1-x\),
\[
I=\int_0^1\frac{(1-x)^2}{x^2+(1-x)^2}\,dx.
\]
Adding gives \(2I=1\), hence \(I=1/2\).
\end{examplebox}

\newpage
\subsubsection{Reciprocal substitution on \((0,\infty)\)}
\begin{extension}
On \((0,\infty)\), the natural symmetry is \(x\leftrightarrow1/x\).  This often
reveals cancellation in logarithmic or palindromic rational integrals.
\end{extension}

\begin{keyformula}[title={Reciprocal Substitution}]
For suitable \(f\),
\[
\int_0^\infty f(x)\,dx
=\int_0^\infty f(1/x)\frac{dx}{x^2}.
\]
\end{keyformula}

\begin{examplebox}[title={Odd Logarithm Under Inversion}]
\[
I=\int_0^\infty\frac{\ln x}{1+x^2}\,dx.
\]
Under \(x\mapsto1/x\),
\[
I=\int_0^\infty\frac{-\ln x}{1+x^2}\,dx=-I,
\]
so \(I=0\).
\end{examplebox}

\newpage
\subsubsection{Homogeneity and scale extraction}
\begin{extension}
If a parameter enters only by scale, remove the scale first.  The remaining
integral is usually dimensionless and may match a standard family.
\end{extension}

\begin{keyformula}[title={Scale-Normalisation Templates}]
\[
\int_0^\infty\frac{x^{m-1}}{(a+x)^n}\,dx
=a^{m-n}\int_0^\infty\frac{t^{m-1}}{(1+t)^n}\,dt,
\qquad x=at,
\]
\[
\int_0^a x^m(a-x)^n\,dx
=a^{m+n+1}\int_0^1 t^m(1-t)^n\,dt,
\qquad x=at.
\]
\end{keyformula}

\begin{examtip}
Scale extraction is especially effective before Beta/Gamma recognition.
\end{examtip}

\newpage
\subsection{Parameter Methods and Differentiation Under the Integral Sign}

\subsubsection{Differentiation under the integral sign}
\begin{extension}
Instead of evaluating one hard integral directly, embed it into a parameter
family, differentiate with respect to the parameter, solve the easier problem,
and then integrate back.
\end{extension}

\begin{keyformula}[title={Parameter Differentiation Template}]
If the interchange is justified, then
\[
F(\lambda)=\int_a^b f(x,\lambda)\,dx
\quad\Rightarrow\quad
F'(\lambda)=\int_a^b \frac{\partial f}{\partial\lambda}(x,\lambda)\,dx.
\]
\end{keyformula}

\begin{examtip}
Choose a parameter whose derivative simplifies the integrand.  Recover the
constant of integration from a convenient value such as \(F(0)\), \(F(1)\), or
a limiting value.
\end{examtip}

\begin{examplebox}[title={Exponential Damping Parameter}]
Let
\[
F(a)=\int_0^\infty\frac{e^{-ax}\sin x}{x}\,dx,
\qquad a>0.
\]
Then
\[
F'(a)=-\int_0^\infty e^{-ax}\sin x\,dx=-\frac{1}{1+a^2}.
\]
Thus \(F(a)=-\arctan a+C\).  Since \(F(a)\to0\) as \(a\to\infty\),
\[
F(a)=\frac\pi2-\arctan a=\arctan\frac1a.
\]
\end{examplebox}

\newpage
\subsubsection{Parameter generation of logarithmic integrals}
\begin{extension}
Powers of \(\ln x\) can be generated by differentiating powers \(x^a\) with
respect to the exponent.
\end{extension}

\begin{keyformula}[title={Logarithmic Integral Generator}]
\[
\frac{\partial^n}{\partial a^n}x^a=x^a(\ln x)^n,
\]
so, for \(a>0\),
\[
\int_0^1 x^{a-1}(\ln x)^n\,dx
=\frac{d^n}{da^n}\left(\frac1a\right)
=\frac{(-1)^n n!}{a^{n+1}}.
\]
\end{keyformula}

\begin{examplebox}[title={Special Values}]
At \(a=1\),
\[
\int_0^1(\ln x)^n\,dx=(-1)^n n!,
\]
so
\[
\int_0^1\ln x\,dx=-1,
\qquad
\int_0^1(\ln x)^2\,dx=2.
\]
\end{examplebox}

\newpage
\subsubsection{Frullani's identity}
\begin{extension}
Frullani's identity is optional enrichment.  It is useful for improper integrals
where a difference cancels the singularity at zero and only a scale ratio
remains.
\end{extension}

\begin{keyformula}[title={Frullani Identity}]
If the relevant limits exist and the integral converges, then for \(a,b>0\),
\[
\int_0^\infty\frac{f(ax)-f(bx)}{x}\,dx
=\bigl(f(0)-f(\infty)\bigr)\ln\frac{b}{a}.
\]
\end{keyformula}

\begin{examplebox}[title={Exponential Frullani Integral}]
For \(f(t)=e^{-t}\),
\[
\int_0^\infty\frac{e^{-ax}-e^{-bx}}{x}\,dx=\ln\frac{b}{a}.
\]
\end{examplebox}

\newpage
\subsection{Reduction Formulas and Special Integral Families}

\subsubsection{Reduction formulas via recurrence design}
\begin{extension}
When one integration-by-parts step returns the same integral family with a lower
exponent, write a recurrence instead of treating each exponent separately.
\end{extension}

\begin{keyformula}[title={Wallis Sine-Power Recurrence}]
\[
I_n:=\int_0^{\pi/2}\sin^n x\,dx,
\qquad
I_n=\frac{n-1}{n}I_{n-2}\quad(n\ge2),
\qquad
I_0=\frac\pi2,
\quad I_1=1.
\]
Hence
\[
I_{2m}=\frac{(2m-1)!!}{(2m)!!}\cdot\frac\pi2,
\qquad
I_{2m+1}=\frac{(2m)!!}{(2m+1)!!}.
\]
\end{keyformula}

\begin{examtip}
A reduction formula solution must include base cases; otherwise the recurrence
has not actually evaluated the integral.
\end{examtip}

\begin{examplebox}[title={Using the Recurrence}]
\[
I_4=\frac34 I_2=\frac34\cdot\frac12 I_0=\frac{3\pi}{16}.
\]
\end{examplebox}

\newpage
\subsubsection{Beta--Gamma bridge}
\begin{extension}
Many difficult-looking definite integrals are disguised Beta or Gamma integrals.
This is enrichment beyond routine MH1101 integration, but it is valuable for
recognising closed forms.
\end{extension}

\begin{keyformula}[title={Beta and Gamma Functions}]
\[
\Gamma(s)=\int_0^\infty t^{s-1}e^{-t}\,dt\quad(s>0),
\qquad
B(p,q)=\int_0^1x^{p-1}(1-x)^{q-1}\,dx\quad(p,q>0),
\]
\[
B(p,q)=\frac{\Gamma(p)\Gamma(q)}{\Gamma(p+q)}.
\]
\end{keyformula}

\begin{keyformula}[title={Common Beta Forms}]
\[
\int_0^\infty\frac{x^{p-1}}{(1+x)^{p+q}}\,dx=B(p,q),
\qquad
\int_0^{\pi/2}\sin^{m-1}x\cos^{n-1}x\,dx
=\frac12 B\!\left(\frac m2,\frac n2\right).
\]
\end{keyformula}

\begin{examplebox}[title={Beta Integral Recognition}]
\[
\int_0^\infty\frac{\sqrt{x}}{(1+x)^3}\,dx
=B\!\left(\frac32,\frac32\right)
=\frac{\Gamma(3/2)^2}{\Gamma(3)}=\frac{\pi}{8}.
\]
\end{examplebox}

\newpage
\subsubsection{Gamma reflection and the \(\pi/\sin(\pi p)\) integral}
\begin{extension}
This is an optional closed-form shortcut for a common improper rational-power
integral.
\end{extension}

\begin{keyformula}[title={Reflection Formula Consequence}]
For \(0<p<1\),
\[
\Gamma(p)\Gamma(1-p)=\frac{\pi}{\sin(\pi p)},
\qquad
\int_0^\infty\frac{x^{p-1}}{1+x}\,dx=\frac{\pi}{\sin(\pi p)}.
\]
\end{keyformula}

\begin{examplebox}[title={Special Case}]
With \(p=1/2\),
\[
\int_0^\infty\frac{dx}{\sqrt{x}(1+x)}=\pi.
\]
\end{examplebox}

\newpage
\subsubsection{Gaussian square trick}
\begin{extension}
The Gaussian integral is a model example where a one-variable integral is best
evaluated by squaring it and using a two-variable polar-coordinate argument.
\end{extension}

\begin{keyformula}[title={Gaussian Integrals}]
\[
\int_{-\infty}^{\infty}e^{-x^2}\,dx=\sqrt\pi,
\qquad
\int_0^\infty e^{-ax^2}\,dx=\frac12\sqrt{\frac\pi a}\quad(a>0),
\]
\[
\int_0^\infty x^m e^{-ax^2}\,dx
=\frac12 a^{-(m+1)/2}\Gamma\!\left(\frac{m+1}{2}\right)
\quad(a>0,\ m>-1).
\]
\end{keyformula}

\begin{examplebox}[title={Square Trick}]
If \(I=\int_{-\infty}^{\infty}e^{-x^2}\,dx\), then
\[
I^2=\iint_{\mathbb R^2}e^{-(x^2+y^2)}\,dA
=\int_0^{2\pi}\int_0^\infty e^{-r^2}r\,dr\,d\theta
=\pi,
\]
so \(I=\sqrt\pi\).
\end{examplebox}

\newpage
\subsubsection{Weighted logarithm antiderivative families}
\begin{extension}
Weighted logarithm integrals can be reduced to exponential-polynomial integrals
by \(x=e^u\), or generated recursively by integration by parts.
\end{extension}

\begin{keyformula}[title={Weighted Logarithm Family}]
For \(n\in\mathbb Z_{\ge0}\),
\[
\int\frac{(\ln x)^n}{x}\,dx=\frac{(\ln x)^{n+1}}{n+1}+C.
\]
For \(m\ne1\),
\[
\int x^{-m}(\ln x)^n\,dx
=-x^{1-m}\sum_{k=0}^{n}
\frac{n!}{(n-k)!(m-1)^{k+1}}(\ln x)^{n-k}+C.
\]
\end{keyformula}

\begin{examplebox}[title={First Logarithmic Cases}]
\[
\int\frac{\ln x}{x^2}\,dx=-\frac{\ln x+1}{x}+C,
\qquad
\int\frac{\ln x}{x^3}\,dx=-\frac{2\ln x+1}{4x^2}+C.
\]
\end{examplebox}

\newpage
\subsection{Series and Taylor Techniques}

\subsubsection{Choosing a convergence test}
\begin{extension}
Convergence testing is a recognition problem.  The fastest solution usually
comes from matching the term structure to the correct test rather than trying
several tests mechanically.
\end{extension}

\begin{examtip}
Use this order:
\begin{enumerate}[leftmargin=*]
\item Check \(a_n\to0\).
\item Check for geometric or \(p\)-series form.
\item Use ratio test for factorials/exponentials.
\item Use root test for \(n\)-th powers.
\item Use comparison or limit comparison for rational/asymptotic terms.
\item For alternating series, check both alternating convergence and absolute convergence.
\end{enumerate}
\end{examtip}

\begin{commonmistake}
Alternating convergence does not imply absolute convergence.  Always test
\(\sum |a_n|\) separately when asked for absolute or conditional convergence.
\end{commonmistake}

\newpage
\subsubsection{Endpoint testing for power series}
\begin{extension}
The radius of convergence gives only the open interval.  Endpoint behaviour is
a separate numerical-series problem.
\end{extension}

\begin{examtip}
For \(\sum c_n(x-a)^n\):
\begin{enumerate}[leftmargin=*]
\item Use ratio/root test to find \(|x-a|<R\).
\item Test \(x=a-R\) as an ordinary series.
\item Test \(x=a+R\) as an ordinary series.
\item Write the final interval using brackets only for convergent endpoints.
\end{enumerate}
\end{examtip}

\begin{commonmistake}
Termwise differentiation and integration preserve the radius but may change
endpoint convergence.  Recheck endpoints after the operation.
\end{commonmistake}

\newpage
\subsubsection{Deriving new series from known series}
\begin{extension}
Known Maclaurin series can generate many new series by substitution,
differentiation, and integration inside the radius of convergence.
\end{extension}

\begin{examtip}
Start from a standard series such as
\[
\frac1{1-x}=\sum_{n=0}^\infty x^n,
\quad
\ln(1+x)=\sum_{n=1}^\infty(-1)^{n+1}\frac{x^n}{n},
\quad
\arctan x=\sum_{n=0}^\infty(-1)^n\frac{x^{2n+1}}{2n+1}.
\]
Then replace \(x\), differentiate, or integrate termwise as needed.
\end{examtip}

\begin{examplebox}[title={Series by Integration}]
Starting from
\[
\frac1{1+x^2}=\sum_{n=0}^\infty(-1)^n x^{2n}\quad(|x|<1),
\]
integrate termwise to get
\[
\arctan x=\sum_{n=0}^\infty(-1)^n\frac{x^{2n+1}}{2n+1}.
\]
\end{examplebox}

\newpage
\subsubsection{Taylor approximation with controlled error}
\begin{extension}
A Taylor polynomial is useful only when paired with a point, degree, and error
control.  The approximation is local around the expansion point.
\end{extension}

\begin{examtip}
To guarantee error below \(\varepsilon\), find \(M\) with
\(|f^{(n+1)}|\le M\) on the interval between \(a\) and \(x\), then solve
\[
\frac{M}{(n+1)!}|x-a|^{n+1}\le\varepsilon.
\]
\end{examtip}

\begin{examplebox}[title={Small Exponential Approximation}]
Using \(e^x=1+x+x^2/2+\cdots\),
\[
T_2(0.1)=1+0.1+\frac{0.1^2}{2}=1.105.
\]
Since \(|e^\xi|\le e^{0.1}<2\) for \(0\le\xi\le0.1\),
\[
|R_2(0.1)|\le\frac{2(0.1)^3}{3!}<0.00034.
\]
\end{examplebox}

\newpage
\subsection{Nonstandard Geometry and General Exam Strategy}

\subsubsection{Pappus centroid theorems}
\begin{extension}
Pappus' theorems are optional shortcuts for volumes and surface areas of
revolution.  They are useful when the centroid of the generating region or
curve is obvious.
\end{extension}

\begin{keyformula}[title={Pappus Theorems}]
If a plane region of area \(A\) is revolved about an external axis and its
centroid is distance \(d\) from the axis, then
\[
V=2\pi dA.
\]
If a plane curve of length \(L\) is revolved about an external axis and its
centroid is distance \(d\) from the axis, then
\[
S=2\pi dL.
\]
For a torus generated by revolving a circle of radius \(r\) whose centre is
at distance \(R>r\) from the axis,
\[
V=2\pi^2Rr^2,
\qquad
S=4\pi^2Rr.
\]
\end{keyformula}

\begin{commonmistake}
The rotation axis must not intersect the generating region or curve in the
standard Pappus setup.
\end{commonmistake}

\newpage
\subsubsection{General exam method-selection checklist}
\begin{extension}
A difficult calculus question often becomes manageable after identifying the
problem type and the most natural representation.
\end{extension}

\begin{examtip}
Use the following decision checklist:
\begin{itemize}[leftmargin=*]
\item Before differentiating, check whether algebraic simplification or a known limit applies.
\item Before integrating, check whether the question asks for setup, exact value, approximation, or convergence only.
\item Before using a powerful theorem, check its hypotheses: continuity, differentiability, positivity, monotonicity, or convergence.
\item In application problems, track signs, units, domains, and endpoints.
\item In series problems, distinguish value, convergence, radius, interval, and error estimate.
\end{itemize}
\end{examtip}

\begin{commonmistake}
Do not use a method merely because it is powerful.  L'H\^opital, Taylor series,
comparison tests, and parameter methods all require hypotheses and often have
simpler alternatives.
\end{commonmistake}

\begin{examplebox}[title={Method Choice Examples}]
\begin{itemize}[leftmargin=*]
\item \(\sqrt{x^2+x}-x\): rationalise before differentiating.
\item \(\int x e^{x}\,dx\): integration by parts.
\item \(\sum n!/n^n\): ratio or root test.
\item \(\sum c_n(x-a)^n\): radius first, endpoints second.
\item Area between curves: split at intersections before integrating.
\end{itemize}
\end{examplebox}
